\section{局限性与未来展望}
\label{zh:sec:limitations}

尽管 MIKU 显著增强了多障碍物交通环境下的交互感知约束生成能力，本文方法仍存在若干实际局限，为未来研究指明了方向。

首先，在预测模型与道路几何方面，本文目前的实证评估主要基于结构化道路环境及带有有界噪声包络的恒定速度预测模型。真实交通场景中包含复杂的道路曲率、拓扑交叉口以及周围人类驾驶员的多模态行为不确定性。这一边界在预测误差呈现系统性而非有界性时可能高估碰撞率的降低幅度。后续研究将引入在自然驾驶数据集上训练的校准多模态轨迹预测模型，例如 Trajectron++ 或 MTR，并将区间几何抽象扩展至曲率超过 0.05\,rad/m 的弯道拓扑。

其次，在验证层级方面，当前的实验验证聚焦于大规模数值仿真与 Apollo 11.0 架构内的软件在环测试。后续工作将补充硬件在环时延、传感器处理延迟以及极端天气或退化路面下的底盘动力学，并开展封闭试验场实车测试与通行顺序动态重排下的完备性分析。

\section{结论}
\label{zh:sec:conclusion}

本文提出了面向动态多障碍交通路径—速度规划的 MIKU 交互感知约束重构框架。通过引入到达时间依赖投影、连续群组级侧向划分、威胁感知安全裕度以及时间同伦图，MIKU 在严格保持下游路径和速度二次规划凸优化命题结构完全不变的前提下，系统性地重构了缺失的交互时序与空间拓扑约束接口。此外，固定截面最大间隙定理保证了最优侧向分配可通过 $O(k\log k)$ 扫描高效求解，消除了传统多障碍物分配中的指数级组合爆炸。

在覆盖七类典型交通构型的 3,500 个配对仿真场景中，实证评估表明 MIKU 将无碰撞到达目标率从 62.7\% 显著提升至 77.5\%，相对基线提升 14.8 个百分点，并将碰撞率从 2.2\% 减半至 1.1\%。性能增益在极具挑战的狭窄通道提升 65.4 个百分点，在延迟交叉场景提升 26.8 个百分点，直接证实了重构后的约束接口能够有效消除虚假死锁。消融实验证实群组级侧向协调与威胁感知裕度构筑了空间可行走廊的基础，而时间同伦图则有效解析了因果通行时序。算法完整规划周期的中位时延仅为 9.98 ms，充分满足自动驾驶典型控制周期的硬实时要求。

在 Apollo 11.0 架构内的软件在环验证进一步证实，MIKU 无需改动原生二次规划求解器即可无缝接入量产级自动驾驶系统，在动态超车、窄通道与车道绕行等代表性工况下表现出优异的兼容性与避障平滑性。闭环滚动重规划实验亦验证了多周期动态行驶下的轨迹一致性与鲁棒性。后续研究将把该框架拓展至校准的多模态轨迹预测、复杂立体弯道道路网络以及全系统实车在环测试平台。
